% ungetestet
%
% Firmware: 0.6~stable20141018
% Gerät: TP-Link TL-WR841N

\documentclass{article}

\usepackage[utf8]{inputenc}	%Zeichencodierung Text
\usepackage[T1]{fontenc}
\newcommand{\changefont}[3]{
\fontfamily{#1} \fontseries{#2} \fontshape{#3} \selectfont}
\usepackage{textcomp}
\usepackage[ngerman]{babel}	%Wortdefinitionen
\usepackage[a5paper,%
 left=20mm,%
 right=20mm,%
 top=20mm,%
 bottom=20mm]{geometry}
\usepackage{fancyhdr}
\usepackage{graphicx}

\pagestyle{fancy}
\fancyhf{}
\rhead{www.freifunk-muenchen.de}
\lhead{\includegraphics[keepaspectratio=true, scale=0.05]{Ffmuc-logo-250.png} Freifunk München}
\rfoot{Mehr Informationen zur Einrichtung unter www.freifunk-muenchen.de}

\renewcommand{\headrulewidth}{1pt}
\renewcommand{\footrulewidth}{1pt}

\renewcommand{\thefootnote}{\roman{footnote}}

\begin{document}

\changefont{phv}{m}{n}



\section*{Kurzanleitung zur Inbetriebnahme Deines Freifunk Routers}
Am Beispiel eines \textbf{TP-Link TL-WR841N} (Stand: 05.11.2014)

\begin{enumerate}
\item Schließe den Freifunk-Router mit dem mitgelieferten Netzteil an das 
    Stromnetz an und schalte ihn ein.

\item Stecke das mitgelieferte LAN-Kabel mit dem einen Ende in die \textbf{LAN}-Buchse deines Computers und mit dem anderen in die blaue \textbf{WAN}-Buchse des Freifunk-Routers.

\item Rufe im Browser Deines Computers \textbf{http://192.168.1.1} auf. \label{webinterface}

\item Konfiguriere deinen Freifunk-Router. 
\begin{enumerate}
  \item Gib deinem Router einen Namen oder akzeptiere die vorgeschlagene Bezeichnung.

  \item Optional können Geokoordinaten\footnote{Mit Hilfe des Werkzeugs \textit{Koordinaten beim nächsten Klicken anzeigen} der Knotenkarte auf freifunk-muenchen.de ermittelbar.} für den Standort des Freifunk-Routers und Kontaktdaten angegeben werden.

  \item Klicke unten rechts auf \glqq{}Fertig\grqq{}.

  \item Nun wird dir der Schlüssel Deines Nodes angezeigt.
  Klicke auf angegebenen Link um den Schlüssel in unser System einzutragen.
  \footnote{Die Konfigurationsseite des Freifunk-Routers ist nun vorerst nicht mehr erreichbar.Durch 3-5 sekündiges Drücken der Reset-Taste wird der Konfigurationsmodus des Routers wieder aktiviert. Anschließend kann wie in Schritt \ref{webinterface} fortgefahren werden.}
\end{enumerate}

\item Verbinde nun die blaue \textbf{WAN}-Buchse des Freifunk-Routers mit einer \textbf{LAN}-Buchse Deines Heimrouters, der Verbindung zum Internet hat. \footnote{Der Freifunk-Router kann sich so mit Hilfe des Mesh-VPN über das Internet mit dem übrigen Freifunk München Intranet verbinden.}

\item Fertig!
\end{enumerate}

\end{document}
